\begin{parts}

\part
می‌خواهیم بررسی کنیم آیا برای سیاست \(\pi_3\) داریم:

\[
V_3 = \frac{1}{2}V_1 + \frac{1}{2}V_2
\]

پاسخ منفی است. برای رد این گزاره یک مثال نقض می‌زنیم.

یک \lr{MDP} با یک حالت غیرپایانی \(s\) و یک حالت پایانی \(t\) در نظر می‌گیریم.
در حالت \(s\) دو کنش \(a\) و \(b\) داریم:

\[
a: \quad s \to s,\qquad r=0
\]

\[
b: \quad s \to t,\qquad r=1
\]

حالت \(t\) پایانی است و ارزش آن صفر است. ضریب تخفیف را هم برابر
\(\gamma = 0.9\) در نظر می‌گیریم.

سیاست اول \(\pi_1\) همیشه کنش \(a\) را انتخاب می‌کند. پس عامل همیشه در حالت
\(s\) می‌ماند و هیچ پاداشی نمی‌گیرد. بنابراین:

\[
V_1(s)=0
\]

سیاست دوم \(\pi_2\) همیشه کنش \(b\) را انتخاب می‌کند. پس عامل همان اول پاداش
\(1\) می‌گیرد و وارد حالت پایانی می‌شود. بنابراین:

\[
V_2(s)=1
\]

اگر رابطه‌ی داده‌شده درست بود، باید می‌داشتیم:

\[
\frac{1}{2}V_1(s)+\frac{1}{2}V_2(s)
=
\frac{1}{2}(0)+\frac{1}{2}(1)
=
\frac{1}{2}
\]

حالا سیاست \(\pi_3\) را بررسی می‌کنیم. این سیاست در هر بار حضور در حالت \(s\)،
با احتمال \(\frac{1}{2}\) کنش \(a\) و با احتمال \(\frac{1}{2}\) کنش \(b\) را انتخاب می‌کند.

پس معادله بلمن برای \(V_3(s)\) به صورت زیر است:

\[
V_3(s)
=
\frac{1}{2}\bigl(0+\gamma V_3(s)\bigr)
+
\frac{1}{2}\bigl(1+\gamma V_3(t)\bigr)
\]

چون \(t\) حالت پایانی است، داریم:

\[
V_3(t)=0
\]

پس:

\[
V_3(s)
=
\frac{1}{2}\gamma V_3(s)
+
\frac{1}{2}
\]

با جایگذاری \(\gamma=0.9\):

\[
V_3(s)
=
0.45V_3(s)+0.5
\]

\[
\begin{aligned}
V_3(s)-0.45V_3(s) &= 0.5 \\
0.55V_3(s) &= 0.5 \\
V_3(s) &= \frac{0.5}{0.55} \\
V_3(s) &= \frac{10}{11}
\end{aligned}
\]

اما از طرف دیگر داشتیم:

\[
\frac{1}{2}V_1(s)+\frac{1}{2}V_2(s)=\frac{1}{2}
\]

پس:

\[
V_3(s)=\frac{10}{11}
\neq
\frac{1}{2}
=
\frac{1}{2}V_1(s)+\frac{1}{2}V_2(s)
\]

بنابراین نمی‌توان گفت تابع ارزش سیاست مخلوط، حتماً برابر با میانگین تابع ارزش
دو سیاست اولیه است. دلیلش هم این است که سیاست جدید در هر گام دوباره تصمیم
تصادفی می‌گیرد و این باعث می‌شود مسیرهایی ساخته شوند که رفتارشان صرفاً میانگین
دو سیاست قبلی نیست.

\end{parts}